\scanpage{135}
lemme 1 et l'hyp. de récurrence,
O.K. pour \(R^ni_*(Q^*)\).
Quant à \(R^ni_*(Rj_*(F))\)
c'est aussi \(R^n(ij)_*(F)\), ce qui nous ramène au cas où \(U\) est \uline{régulier}.
Alors par \uuline{résolution}, on a \(i=fi'\),
avec \(f\) projectif, \(X'\) régulier, \(i'\) une immersion ouverte à complémentaire
un diviseur à croisements normaux.

\diagram{\[
\begin{tikzcd}[column sep=small]
&X'\arrow[d,"f"]\\
U\arrow[ur,hook,"i'"]&X\arrow[l,"i"]
\end{tikzcd}
\]}

Alors
\[
R^\alpha fi'_*(F)\Longleftarrow R^pf_*(R^qi'_*(F))
\]
et cela nous ramène \inserted{par le lemme 3} au cas où
\deleted{\(\unsure{E}\quad\unsure{X}\)}
\(X-U\) est un diviseur à croisements normaux.
\uuline{Lorsque} \(F=Q_l\), on gagne alors par le calcul explicite,
\inserted{(th. de \uuline{pureté}).}
On va ramener le cas général à celui-ci, en complétant \(V/\deleted{X}\) en
\deleted{\(\unsure{f}\)}
\(g:Y\mathbin{\text{---}}X\) morphisme projectif (pas lisse en général),
où \(F\) est un facteur direct de \(R^nf_*(Q_l)\),
\inserted{en supposant dim rel. \(g\le d'\)}.
On peut supposer
\[
F=R^nf_*(Q_l)=R^ng_*(Q_l)|U.
\]
\diagram{\[
\begin{tikzcd}
V\arrow[r,hook,no head,"j"]\arrow[d,"f"']&
\unsure{\widehat Y|Y}\arrow[d,"g"]\\
U\arrow[r,hook,no head,"i"']&X
\end{tikzcd}
\]}
Par des dévissages faciles, utilisant lemme 3,
on est ramené au cas où \(d'=1\).
\inserted{et \(f\) \unsure{à|:} fibres géom. conn.}
\margin{\uline{à vérifier}}
Alors la suite spectrale pour \(\unsure{if|ig}\) est
% archive foliation: 134
